% Part: turing-machines
% Chapter: machines-computations
% Section: combining-machines

\documentclass[../../../include/open-logic-section]{subfiles}

\begin{document}

\olfileid{tur}{mac}{cmb}
\olsection{ٹیورنگ مشیناں نوں جوڑنا}

\begin{explain}
ٹیورنگ مشیناں دیاں جیہڑیاں مثالاں اسی ہن تک ویکھیاں نیں،
اوہ خاصیاں سادیاں سن۔ پر حقیقت وچ جیہڑا مسئلہ وی
کسے جدید پروگرامنگ زبان نال حل کیتا جا سکدا اے،
اوہ ٹیورنگ مشیناں نال وی حل کیتا جا سکدا اے۔
ہور پیچیدہ ٹیورنگ مشیناں بنانے لئی ایہ سمجھنا
اہم اے کہ اسی مشیناں نوں آپس وچ جوڑ سکدے آں:
عمل نوں سادے حصیاں وچ ونڈ کے وڈے مسئلے
حل کرن والی مشین بنا سکدے آں۔
جے پیچیدہ مسئلے نوں فطری انداز وچ حصیاں وچ
ونڈیا جا سکدا ہووے، تاں اسی کئی مرحلیاں وچ
ہر حصے لئی سادی مشین بنا کے اوہناں نوں اک
مشین وچ جوڑ سکدے آں جیڑی پورا مسئلہ حل کرے۔
اگلے حصے وچ رُکّن دے مسئلے نال نبڑدیاں ایہ گل
خاص طور اُتے اہم ہووے گا۔

ٹیورنگ مشیناں $M = \tuple{Q, \Sigma, q_0, \delta}$
تے~$M' = \tuple{Q', \Sigma', q_0', \delta'}$
نوں اسی کِویں جوڑئیے؟ ہن $M$~دے رُکّن مگروں
پٹی دی مکمل حالت نوں مشین~$M'$ دے چلن دی
شروع والی مکمل حالت بنا لیندے آں۔
ایہ کم کرن والی اک ٹیورنگ مشین $M \frown M'$
بنان لئی ایہ قدم چکو:
\begin{enumerate}
    \item حالتاں~$Q'$ نوں، جیڑیاں $M'$ دیاں نیں،
    نویں نمبر یا ناں دیو، تاں جو $M$
    تے~$M'$ دی کوئی حالت سانجھی نہ ہووے
    ($Q \cap Q' = \emptyset$)۔
    \item $M \frown M'$ دیاں حالتاں $Q \cup Q'$ نیں۔
    \item پٹی دے نشاناں دا مجموعہ $\Sigma \cup \Sigma'$ اے۔
    \item ابتدائی حالت~$q_0$ اے۔
    \item حالتی تبدیلی والا فنکشن $\delta''$
    ایہ اے:
    \[\delta''(q,\sigma) =
    \begin{cases}
      \delta(q,\sigma) & \text{جے $q \in Q$ ہووے تے $\delta(q,\sigma)$ تعریف شدہ ہووے}\\
      \delta'(q,\sigma) & \text{جے $q \in Q'$ ہووے}\\
      \tuple{q_0', \sigma, \TMstay} & \text{جے $q \in Q$ ہووے تے
      $\delta(q,\sigma)$ غیر تعریف شدہ ہووے}
    \end{cases}\]
\end{enumerate}
جوڑی ہوئی مشین وچ $M$ دا حکم ورتیا جاندا اے
جدوں حالت $q \in Q$ ہووے تے اوس ویلے کوئی حکم تعریف شدہ ہووے؛
$M'$ دے حکم حالت~$q \in Q'$ وچ ورتے جاندے نیں۔
جے حالت $q \in Q$ ہووے تے سر اوہ
نشان~$\sigma$ پڑھدا ہووے جیہدے لئی $M$
دی حالتی تبدیلی تعریف شدہ نہیں
(یعنی $M$ ایتھے رُک جاندی)،
تاں مشین~$M'$ دی ابتدائی حالت وچ
چلی جاندی اے۔ پٹی دے نشان
تے سر دی جگہ اوہو رہندے نیں۔

یاد رکھو: جے مشین~$M$ پابند نہ ہووے،
تاں $M$~دے رُکّن ویلے سر دی جگہ دا
سانوں پتا نہیں ہوندا۔ ایس لئی
$M$~دی آخری مکمل حالت وچ سر دا
خانہ~$1$ پڑھنا ضروری نہیں۔
مشیناں جوڑدیاں ایہ گل ذہن وچ رکھنی چاہیدی اے۔
\end{explain}

\begin{ex}
\emph{مشیناں نوں جوڑنا:} اسی ایہو جیہی اک مشین بناواں گے
جیڑی $\TMstroke$ دیاں دو لڑیاں،
اک دی لمبائی $n$ تے دوجی دی $m$،
انپٹ وچ لے کے رُکے تاں پٹی اُتے
$2(m+n)$ $\TMstroke$ دی اکّو لڑی ہووے۔
ایہ بناؤن لئی اسی دو پہلاں توں
جانیاں مشیناں جوڑ سکدے آں:
جمع والی مشین تے دوگنا کرن والی مشین۔
پہلاں جمع والی مشین دا حالتی خاکہ بناندے آں۔
\[
\begin{tikzpicture}[->,>=stealth',shorten >=1pt,auto,node distance=2.8cm,
                    semithick]
  \tikzstyle{every state}=[fill=none,draw=black,text=black]

  \node[initial,state] (A)              {$q_0$};
  \node[state]         (B) [right of=A] {$q_1$};
  \node[state]         (C) [right of=B] {$q_2$};

  \path (A) edge node {\TMtrans{\TMblank}{\TMstroke}{\TMstay}} (B)
            edge [loop above] node {\TMtrans{\TMstroke}{\TMstroke}{\TMright}} (B)
        (B) edge [loop above] node {\TMtrans{\TMstroke}{\TMstroke}{\TMright}} (B)
            edge node {\TMtrans{\TMblank}{\TMblank}{\TMleft}} (C)
        (C) edge [loop above] node {\TMtrans{\TMstroke}{\TMblank}{\TMstay}} (C);
\end{tikzpicture}
\]
حالت~$q_2$ وچ رُکّن دی بجائے
اسی عمل جاری رکھنا چاہندے آں
تاں جو حاصل نوں دوگنا کیتا جا سکے۔
یاد کرو کہ دوگنا کرن والی مشین
انپٹ دی پہلی لکیر مٹا کے وکھرے
آؤٹ پٹ وچ دو لکیراں لکھدی اے۔
اک حکم ہور شامل کریے تاں جو
سر جمع والی مشین دے آؤٹ پٹ
دی پہلی لکیر پڑھ رہیا ہووے۔
\[
\begin{tikzpicture}[->,>=stealth',shorten >=1pt,auto,node distance=2.8cm,
                    semithick]
  \tikzstyle{every state}=[fill=none,draw=black,text=black]

  \node[initial,state] (A)              {$q_0$};
  \node[state]         (B) [right of=A] {$q_1$};
  \node[state]         (C) [right of=B] {$q_2$};
  \node[state]         (D) [below left of=C] {$q_3$};
  \node[state]         (E) [below left of=D] {$q_4$};

  \path (A) edge node {\TMtrans{\TMblank}{\TMstroke}{\TMstay}} (B)
            edge [loop above] node {\TMtrans{\TMstroke}{\TMstroke}{\TMright}} (B)
        (B) edge [loop above] node {\TMtrans{\TMstroke}{\TMstroke}{\TMright}} (B)
            edge node {\TMtrans{\TMblank}{\TMblank}{\TMleft}} (C)
        (C) edge node {\TMtrans{\TMstroke}{\TMblank}{\TMleft}} (D)
        (D) edge [loop left] node {\TMtrans{\TMstroke}{\TMstroke}{\TMleft}} (D)
            edge node[left, xshift=-2mm] {\TMtrans{\TMendtape}{\TMendtape}{\TMright}} (E);
\end{tikzpicture}
\]
ہن انپٹ نوں دوگنا کرنا سوکھا اے:
دوگنا کرن والی مشین نوں حالت~$q_4$
نال جوڑ دیو۔ ایس لئی اوس مشین دیاں
حالتیں نویں ناں نال لکھنیاں پین گیاں
تاں جو اوہ~$q_4$ توں شروع ہوون،
نہ کہ~$q_0$ توں؛ ایس طرح دو
ابتدائی حالتیں نہیں بنن گیاں۔
آخری خاکہ \olref{fig:combined} ورگا ہووے گا۔
\begin{figure}
\[
\begin{tikzpicture}[->,>=stealth',shorten >=1pt,auto,node distance=2.8cm,
                    semithick]
  \tikzstyle{every state}=[fill=none,draw=black,text=black]
  \node[initial,state] (A)              {$q_0$};
  \node[state]         (B) [right of=A] {$q_1$};
  \node[state]         (C) [right of=B] {$q_2$};
  \node[state]         (D) [below left of=C] {$q_3$};
  \node[state]         (E) [below left of=D] {$q_4$};
  \node[state]         (2) [right of=E] {$q_5$};
  \node[state]         (3) [right of=2] {$q_6$};
  \node[state]         (4) [below of=3] {$q_7$};
  \node[state]         (5) [left of=4]  {$q_8$};
  \node[state]         (6) [left of=5]  {$q_9$};

  \path (A) edge node {\TMtrans{\TMblank}{\TMstroke}{\TMstay}} (B)
            edge [loop above] node {\TMtrans{\TMstroke}{\TMstroke}{\TMright}} (B)
        (B) edge [loop above] node {\TMtrans{\TMstroke}{\TMstroke}{\TMright}} (B)
            edge node {\TMtrans{\TMblank}{\TMblank}{\TMleft}} (C)
        (C) edge node {\TMtrans{\TMstroke}{\TMblank}{\TMleft}} (D)
        (D) edge [loop left] node {\TMtrans{\TMstroke}{\TMstroke}{\TMleft}} (D)
            edge node[left, xshift=-2mm] {\TMtrans{\TMendtape}{\TMendtape}{\TMright}} (E)
    (E) edge node {\TMtrans{\TMstroke}{\TMblank}{\TMright}} (2)
    (2) edge [loop above] node {\TMtrans{\TMstroke}{\TMstroke}{\TMright}} (2)
      edge node {\TMtrans{\TMblank}{\TMblank}{\TMright}} (3)
    (3) edge [loop above] node {\TMtrans{\TMstroke}{\TMstroke}{\TMright}} (3)
        edge node {\TMtrans{\TMblank}{\TMstroke}{\TMright}} (4)
    (4) edge [loop below] node {\TMtrans{\TMblank}{\TMstroke}{\TMleft}} (4)
        edge node {\TMtrans{\TMstroke}{\TMstroke}{\TMleft}} (5)
    (5) edge [loop below]  node {\TMtrans{\TMstroke}{\TMstroke}{\TMleft}} (5)
        edge              node {\TMtrans{\TMblank}{\TMblank}{\TMleft}} (6)
    (6) edge [loop below] node {\TMtrans{\TMstroke}{\TMstroke}{\TMleft}} (6)
        edge              node {\TMtrans{\TMblank}{\TMblank}{\TMright}} (E);
\end{tikzpicture}
\]
\caption{جمع تے دوگنا کرن والیاں مشیناں نوں جوڑنا}
\ollabel{fig:combined}
\end{figure}
\end{ex}

\begin{prop}
    جے $M$ تے $M'$ پابند مشیناں
    بالترتیب فنکشن $f\colon
    \Nat^k \to \Nat$ تے $f'\colon \Nat \to \Nat$
    دا حساب کردیاں نیں، تاں
    $M \frown M'$ وی پابند اے
    تے~$\comp{f}{f'}$ دا حساب کردی اے۔
\end{prop}

\begin{proof}
    چونکہ $M$ پابند اے، اس لئی
    جدوں اوہ آؤٹ پٹ~$f(n_1,\dots,n_k) = m$
    دے کے رُکدی اے تاں سر خانہ~$1$
    پڑھ رہیا ہوندا اے۔ ہن جے اسی
    مشین~$M'$ دی ابتدائی حالت وچ
    داخل ہوئیے، تاں ایہ مشین آؤٹ پٹ
    $f'(m)$ دے کے رُکے گی تے سر فیر
    خانہ~$1$ پڑھدا ہووے گا۔
    \olref[dis]{defn:disciplined} دیاں
    باقی شرطاں وی پوری ہوندیاں نیں۔
\end{proof}

\begin{prob}
    فنکشن $f(x) = x+2$ دا حساب کرن
    والی پابند ٹیورنگ مشین دیو:
    \cref{tur:mac:dis:prob:disc-succ} والی
    مشین~$M$ لے کے $M \frown M$ بناؤ۔
\end{prob}
\end{document}
